\IfFileExists{stacks-project.cls}{%
\documentclass{stacks-project}
}{%
\documentclass{amsart}
}

% For dealing with references we use the comment environment
\usepackage{verbatim}
\newenvironment{reference}{\comment}{\endcomment}
%\newenvironment{reference}{}{}
\newenvironment{slogan}{\comment}{\endcomment}
\newenvironment{history}{\comment}{\endcomment}

% For commutative diagrams we use Xy-pic
\usepackage[all]{xy}

% We use 2cell for 2-commutative diagrams.
\xyoption{2cell}
\UseAllTwocells

% We use multicol for the list of chapters between chapters
\usepackage{multicol}

% This is generally recommended for better output
\usepackage{lmodern}
\usepackage[T1]{fontenc}

% For cross-file-references
\usepackage{xr-hyper}

% Package for hypertext links:
\usepackage{hyperref}

% For any local file, say "hello.tex" you want to link to please
% use \externaldocument[hello-]{hello}
\externaldocument[introduction-]{introduction}
\externaldocument[conventions-]{conventions}
\externaldocument[sets-]{sets}
\externaldocument[categories-]{categories}
\externaldocument[topology-]{topology}
\externaldocument[sheaves-]{sheaves}
\externaldocument[sites-]{sites}
\externaldocument[stacks-]{stacks}
\externaldocument[fields-]{fields}
\externaldocument[algebra-]{algebra}
\externaldocument[brauer-]{brauer}
\externaldocument[homology-]{homology}
\externaldocument[derived-]{derived}
\externaldocument[simplicial-]{simplicial}
\externaldocument[more-algebra-]{more-algebra}
\externaldocument[smoothing-]{smoothing}
\externaldocument[modules-]{modules}
\externaldocument[sites-modules-]{sites-modules}
\externaldocument[injectives-]{injectives}
\externaldocument[cohomology-]{cohomology}
\externaldocument[sites-cohomology-]{sites-cohomology}
\externaldocument[dga-]{dga}
\externaldocument[dpa-]{dpa}
\externaldocument[sdga-]{sdga}
\externaldocument[hypercovering-]{hypercovering}
\externaldocument[schemes-]{schemes}
\externaldocument[constructions-]{constructions}
\externaldocument[properties-]{properties}
\externaldocument[morphisms-]{morphisms}
\externaldocument[coherent-]{coherent}
\externaldocument[divisors-]{divisors}
\externaldocument[limits-]{limits}
\externaldocument[varieties-]{varieties}
\externaldocument[topologies-]{topologies}
\externaldocument[descent-]{descent}
\externaldocument[perfect-]{perfect}
\externaldocument[more-morphisms-]{more-morphisms}
\externaldocument[flat-]{flat}
\externaldocument[groupoids-]{groupoids}
\externaldocument[more-groupoids-]{more-groupoids}
\externaldocument[etale-]{etale}
\externaldocument[chow-]{chow}
\externaldocument[intersection-]{intersection}
\externaldocument[pic-]{pic}
\externaldocument[weil-]{weil}
\externaldocument[adequate-]{adequate}
\externaldocument[dualizing-]{dualizing}
\externaldocument[duality-]{duality}
\externaldocument[discriminant-]{discriminant}
\externaldocument[derham-]{derham}
\externaldocument[local-cohomology-]{local-cohomology}
\externaldocument[algebraization-]{algebraization}
\externaldocument[curves-]{curves}
\externaldocument[resolve-]{resolve}
\externaldocument[models-]{models}
\externaldocument[functors-]{functors}
\externaldocument[equiv-]{equiv}
\externaldocument[pione-]{pione}
\externaldocument[etale-cohomology-]{etale-cohomology}
\externaldocument[proetale-]{proetale}
\externaldocument[relative-cycles-]{relative-cycles}
\externaldocument[more-etale-]{more-etale}
\externaldocument[trace-]{trace}
\externaldocument[crystalline-]{crystalline}
\externaldocument[spaces-]{spaces}
\externaldocument[spaces-properties-]{spaces-properties}
\externaldocument[spaces-morphisms-]{spaces-morphisms}
\externaldocument[decent-spaces-]{decent-spaces}
\externaldocument[spaces-cohomology-]{spaces-cohomology}
\externaldocument[spaces-limits-]{spaces-limits}
\externaldocument[spaces-divisors-]{spaces-divisors}
\externaldocument[spaces-over-fields-]{spaces-over-fields}
\externaldocument[spaces-topologies-]{spaces-topologies}
\externaldocument[spaces-descent-]{spaces-descent}
\externaldocument[spaces-perfect-]{spaces-perfect}
\externaldocument[spaces-more-morphisms-]{spaces-more-morphisms}
\externaldocument[spaces-flat-]{spaces-flat}
\externaldocument[spaces-groupoids-]{spaces-groupoids}
\externaldocument[spaces-more-groupoids-]{spaces-more-groupoids}
\externaldocument[bootstrap-]{bootstrap}
\externaldocument[spaces-pushouts-]{spaces-pushouts}
\externaldocument[spaces-chow-]{spaces-chow}
\externaldocument[groupoids-quotients-]{groupoids-quotients}
\externaldocument[spaces-more-cohomology-]{spaces-more-cohomology}
\externaldocument[spaces-simplicial-]{spaces-simplicial}
\externaldocument[spaces-duality-]{spaces-duality}
\externaldocument[formal-spaces-]{formal-spaces}
\externaldocument[restricted-]{restricted}
\externaldocument[spaces-resolve-]{spaces-resolve}
\externaldocument[formal-defos-]{formal-defos}
\externaldocument[defos-]{defos}
\externaldocument[cotangent-]{cotangent}
\externaldocument[examples-defos-]{examples-defos}
\externaldocument[algebraic-]{algebraic}
\externaldocument[examples-stacks-]{examples-stacks}
\externaldocument[stacks-sheaves-]{stacks-sheaves}
\externaldocument[criteria-]{criteria}
\externaldocument[artin-]{artin}
\externaldocument[quot-]{quot}
\externaldocument[stacks-properties-]{stacks-properties}
\externaldocument[stacks-morphisms-]{stacks-morphisms}
\externaldocument[stacks-limits-]{stacks-limits}
\externaldocument[stacks-cohomology-]{stacks-cohomology}
\externaldocument[stacks-perfect-]{stacks-perfect}
\externaldocument[stacks-introduction-]{stacks-introduction}
\externaldocument[stacks-more-morphisms-]{stacks-more-morphisms}
\externaldocument[stacks-geometry-]{stacks-geometry}
\externaldocument[moduli-]{moduli}
\externaldocument[moduli-curves-]{moduli-curves}
\externaldocument[examples-]{examples}
\externaldocument[exercises-]{exercises}
\externaldocument[guide-]{guide}
\externaldocument[desirables-]{desirables}
\externaldocument[coding-]{coding}
\externaldocument[obsolete-]{obsolete}
\externaldocument[fdl-]{fdl}
\externaldocument[index-]{index}

% Theorem environments.
%
\theoremstyle{plain}
\newtheorem{theorem}[subsection]{Theorem}
\newtheorem{proposition}[subsection]{Proposition}
\newtheorem{lemma}[subsection]{Lemma}

\theoremstyle{definition}
\newtheorem{definition}[subsection]{Definition}
\newtheorem{example}[subsection]{Example}
\newtheorem{exercise}[subsection]{Exercise}
\newtheorem{situation}[subsection]{Situation}

\theoremstyle{remark}
\newtheorem{remark}[subsection]{Remark}
\newtheorem{remarks}[subsection]{Remarks}

\numberwithin{equation}{subsection}

% Macros
%
\def\lim{\mathop{\mathrm{lim}}\nolimits}
\def\colim{\mathop{\mathrm{colim}}\nolimits}
\def\Spec{\mathop{\mathrm{Spec}}}
\def\Hom{\mathop{\mathrm{Hom}}\nolimits}
\def\Ext{\mathop{\mathrm{Ext}}\nolimits}
\def\SheafHom{\mathop{\mathcal{H}\!\mathit{om}}\nolimits}
\def\SheafExt{\mathop{\mathcal{E}\!\mathit{xt}}\nolimits}
\def\Sch{\mathit{Sch}}
\def\Mor{\mathop{\mathrm{Mor}}\nolimits}
\def\Ob{\mathop{\mathrm{Ob}}\nolimits}
\def\Sh{\mathop{\mathit{Sh}}\nolimits}
\def\NL{\mathop{N\!L}\nolimits}
\def\CH{\mathop{\mathrm{CH}}\nolimits}
\def\proetale{{pro\text{-}\acute{e}tale}}
\def\etale{{\acute{e}tale}}
\def\QCoh{\mathit{QCoh}}
\def\Ker{\mathop{\mathrm{Ker}}}
\def\Im{\mathop{\mathrm{Im}}}
\def\Coker{\mathop{\mathrm{Coker}}}
\def\Coim{\mathop{\mathrm{Coim}}}

% Boxtimes
%
\DeclareMathSymbol{\boxtimes}{\mathbin}{AMSa}{"02}

%
% Macros for moduli stacks/spaces
%
\def\QCohstack{\mathcal{QC}\!\mathit{oh}}
\def\Cohstack{\mathcal{C}\!\mathit{oh}}
\def\Spacesstack{\mathcal{S}\!\mathit{paces}}
\def\Quotfunctor{\mathrm{Quot}}
\def\Hilbfunctor{\mathrm{Hilb}}
\def\Curvesstack{\mathcal{C}\!\mathit{urves}}
\def\Polarizedstack{\mathcal{P}\!\mathit{olarized}}
\def\Complexesstack{\mathcal{C}\!\mathit{omplexes}}
% \Pic is the operator that assigns to X its picard group, usage \Pic(X)
% \Picardstack_{X/B} denotes the Picard stack of X over B
% \Picardfunctor_{X/B} denotes the Picard functor of X over B
\def\Pic{\mathop{\mathrm{Pic}}\nolimits}
\def\Picardstack{\mathcal{P}\!\mathit{ic}}
\def\Picardfunctor{\mathrm{Pic}}
\def\Deformationcategory{\mathcal{D}\!\mathit{ef}}
\usepackage{cleveref}
\crefname{originallemma}{original lemma}{original lemmas}
\Crefname{originallemma}{Original lemma}{Original lemmas}
\makeatletter
\newlabel{hypercovering-lemma-hypercovering-morphism-sites}{{0DAY}{source}{}{sourcepointer-1}{}}
\newlabel{hypercovering-lemma-hypercovering-morphism-sites@cref}{{[originallemma][0][]0DAY}{[1][0][]source}{}{}{}}
\makeatother
\crefname{originallemma}{original lemma}{original lemmas}
\Crefname{originallemma}{Original lemma}{Original lemmas}
\makeatletter
\newlabel{hypercovering-lemma-hypercovering-site}{{0DAX}{source}{}{sourcepointer-2}{}}
\newlabel{hypercovering-lemma-hypercovering-site@cref}{{[originallemma][0][]0DAX}{[1][0][]source}{}{}{}}
\makeatother
\crefname{originalreference}{original reference}{original references}
\Crefname{originalreference}{Original reference}{Original references}
\makeatletter
\newlabel{lemma-hypercovering-equivalence-bounded-modules}{{0DA2}{source}{}{sourcepointer-3}{}}
\newlabel{lemma-hypercovering-equivalence-bounded-modules@cref}{{[originalreference][0][]0DA2}{[1][0][]source}{}{}{}}
\makeatother
\crefname{originalreference}{original reference}{original references}
\Crefname{originalreference}{Original reference}{Original references}
\makeatletter
\newlabel{lemma-prepare-bbd-glueing}{{0DCA}{source}{}{sourcepointer-4}{}}
\newlabel{lemma-prepare-bbd-glueing@cref}{{[originalreference][0][]0DCA}{[1][0][]source}{}{}{}}
\makeatother
\crefname{originalreference}{original reference}{original references}
\Crefname{originalreference}{Original reference}{Original references}
\makeatletter
\newlabel{situation-locally-given}{{0D67}{source}{}{sourcepointer-5}{}}
\newlabel{situation-locally-given@cref}{{[originalreference][0][]0D67}{[1][0][]source}{}{}{}}
\makeatother
\title{BBD glueing: locally given derived objects on a basis glue uniquely under negative-Ext and cohomological-bound vanishing}
\author{The Stacks Project Authors}
\date{stacks-project:spaces-simplicial.tex:local-snapshot-20261012; GNU FDL 1.2 or later}
\begin{document}
\maketitle
\noindent\textit{Editorial note.} All delivered mathematical text is the author's own, copied from the pinned Stacks Project chapter file sources/stacks-spaces-simplicial/pinned-original-source.tex; no statement or proof step was rewritten, reordered or removed.\par\medskip
\noindent\textit{Editorial note.} The counted proof is the author's complete argument: the 'small detail omitted' (source line 5705) is the transfer of uniqueness: F is determined by E = a\textasciicircum{}*F through the bounded-modules equivalence (a\textasciicircum{}* is fully faithful with quasi-inverse La\_* on the bounded range), so the unique-isomorphism uniqueness of E gives that of F - an auxiliary step away from the compatibility construction, which is fully argued; this lemma is the engine invoked by card L46160 (the Cech glueing special case).\par\medskip
\noindent\textit{Editorial note.} The label lemma-bbd-glueing names the counted statement, and the proof environment immediately following it in the source is delivered verbatim as its proof.\par\medskip
\noindent\textit{Editorial note.} The counted statement is carried by the author's own lemma environment in the chapter source.\par\medskip
\noindent\textit{Allowed standard prerequisites (editorial).}\begin{itemize}
\item The locally-given situation. Same chapter, Situation \textbackslash{}label\{situation-locally-given\}. Line 5629: the data (B, E\_U, rho\_a) whose solution is to be constructed.
\item Hypercoverings adapted to a full subcategory. Hypercoverings chapter, \textbackslash{}label\{hypercovering-lemma-hypercovering-site\}. Lines 5653-5656: a hypercovering L of D with L\_n = \{U\_\{n,i\}\} inside B; its proof's construction of L\_0, L\_1 is reused at lines 5673-5677.
\item Preparation lemma producing the cartesian object on the total site. Same chapter, \textbackslash{}label\{lemma-prepare-bbd-glueing\}. Lines 5657-5660: the cartesian object E on (C/K)\_total restricting compatibly to the E\_\{U\_\{n,i\}\}; its compatibility statement is reused at lines 5691-5696.
\item Hypercoverings are preserved by morphisms of sites. Hypercoverings chapter, \textbackslash{}label\{hypercovering-lemma-hypercovering-morphism-sites\}. Lines 5662-5663: K = u(L) is a hypercovering.
\item Bounded modules equivalence for hypercoverings. Same chapter, \textbackslash{}label\{lemma-hypercovering-equivalence-bounded-modules\}. Lines 5665-5666: E in D\textasciicircum{}+ is a\textasciicircum{}*F for some F in D\textasciicircum{}+(O\_C).
\end{itemize}\medskip
\ifdefined\hypertarget\hypertarget{sourcepointer-1}{}\fi
\noindent\textit{Source reference.} Original lemma 0DAY: sources/stacks-chow/tags.txt, source lines 16678--16678. Prerequisite pointer invoked at proof line(s) 5662; it names a standard fact of a sibling chapter that the counted argument uses, and every step of the counted proof built on it is retained verbatim in the delivered body.\par\medskip
\ifdefined\hypertarget\hypertarget{sourcepointer-2}{}\fi
\noindent\textit{Source reference.} Original lemma 0DAX: sources/stacks-chow/tags.txt, source lines 16677--16677. Prerequisite pointer invoked at proof line(s) 5654; it names a standard fact of a sibling chapter that the counted argument uses, and every step of the counted proof built on it is retained verbatim in the delivered body.\par\medskip
\ifdefined\hypertarget\hypertarget{sourcepointer-3}{}\fi
\noindent\textit{Source reference.} Original reference 0DA2: sources/stacks-spaces-simplicial/pinned-original-source.tex, source lines 4754--4754. Prerequisite pointer invoked at proof line(s) 5664; it names a standard fact of the same chapter that the counted argument uses, and every step of the counted proof built on it is retained verbatim in the delivered body.\par\medskip
\ifdefined\hypertarget\hypertarget{sourcepointer-4}{}\fi
\noindent\textit{Source reference.} Original reference 0DCA: sources/stacks-spaces-simplicial/pinned-original-source.tex, source lines 5588--5588. Prerequisite pointer invoked at proof line(s) 5657; it names a standard fact of the same chapter that the counted argument uses, and every step of the counted proof built on it is retained verbatim in the delivered body.\par\medskip
\ifdefined\hypertarget\hypertarget{sourcepointer-5}{}\fi
\noindent\textit{Source reference.} Original reference 0D67: sources/stacks-spaces-simplicial/pinned-original-source.tex, source lines 5552--5552. Prerequisite pointer invoked at proof line(s) 5631; it names a standard fact of the same chapter that the counted argument uses, and every step of the counted proof built on it is retained verbatim in the delivered body.\par\medskip
\begin{lemma}[BBD glueing lemma]
\label{lemma-bbd-glueing}
In Situation \ref{situation-locally-given}. Assume
\begin{enumerate}
\item $\mathcal{C}$ has equalizers and fibre products,
\item there is a morphism of sites $f : \mathcal{C} \to \mathcal{D}$
given by a continuous functor $u : \mathcal{D} \to \mathcal{C}$
such that
\begin{enumerate}
\item $\mathcal{D}$ has equalizers and fibre products and $u$
commutes with them,
\item $\mathcal{B}$ is a full subcategory of $\mathcal{D}$
and $u : \mathcal{B} \to \mathcal{C}$ is the restriction of $u$,
\item every object of $\mathcal{D}$ has a covering whose members
are objects of $\mathcal{B}$,
\end{enumerate}
\item for all $U$ in $\mathcal{B}$ all negative self-exts of $E_U$
in $D(\mathcal{O}_{u(U)})$ are zero, and
\item there exists a $t \in \mathbf{Z}$ such that $H^i(E_U) = 0$ for $i < t$
and $U \in \Ob(\mathcal{B})$.
\end{enumerate}
Then there exists a solution unique up to unique isomorphism.
\end{lemma}

\begin{proof}
By Hypercoverings, Lemma \ref{hypercovering-lemma-hypercovering-site}
there exists a hypercovering $L$ for the site $\mathcal{D}$ such that
$L_n = \{U_{n, i}\}_{i \in I_n}$ with $U_{i, n} \in \Ob(\mathcal{B})$.
Set $K = u(L)$. Apply Lemma \ref{lemma-prepare-bbd-glueing}
to get a cartesian object $E$ of $D(\mathcal{O})$ on the site
$(\mathcal{C}/K)_{total}$ restricting to $E_{U_{n, i}}$ on
$\mathcal{C}/u(U_{n, i})$ compatibly.
The assumption on $t$ implies that $E \in D^+(\mathcal{O})$.
By Hypercoverings, Lemma \ref{hypercovering-lemma-hypercovering-morphism-sites}
we see that $K$ is a hypercovering too.
By Lemma \ref{lemma-hypercovering-equivalence-bounded-modules}
we find that $E = a^*F$ for some $F$ in $D^+(\mathcal{O}_\mathcal{C})$.

\medskip\noindent
To prove that $F$ is a solution we will use the construction of
$L_0$ and $L_1$ given in the proof of
Hypercoverings, Lemma \ref{hypercovering-lemma-hypercovering-site}.
(This is a bit inelegant but there does not seem to be a completely
straightforward way around it.)

\medskip\noindent
Namely, we have $I_0 = \Ob(\mathcal{B})$ and so
$L_0 = \{U\}_{U \in \Ob(\mathcal{B})}$.
Hence the isomorphism $a^*F \to E$ restricted to the components
$\mathcal{C}/u(U)$ of $\mathcal{C}/K_0$ defines isomorphisms
$\rho_U : F|_{\mathcal{C}/u(U)} \to E_U$ for $U \in \Ob(\mathcal{B})$
by our choice of $E$.

\medskip\noindent
To prove that $\rho_U$ satisfy the requirement of compatibility
with the maps $\rho_a$ of Situation \ref{situation-locally-given}
we use that $I_1$ contains the set
$$
\Omega =
\{(U, V, W, a, b) \mid U, V, W \in \mathcal{B}, a : U \to V, b : U \to W\}
$$
and that for $i = (U, V, W, a, b)$ in $\Omega$ we have
$U_{1, i} = U$. Moreover, the component maps $f_{\delta^1_0, i}$ and
$f_{\delta^1_1, i}$ of the two morphisms $K_1 \to K_0$ are the morphisms
$$
a : U \to V \quad\text{and}\quad b : U \to V
$$
Hence the compatibility mentioned in
Lemma \ref{lemma-prepare-bbd-glueing} gives that
$$
\rho_a \circ \rho_V|_{\mathcal{C}/u(U)} = \rho_U
\quad\text{and}\quad
\rho_b \circ \rho_W|_{\mathcal{C}/u(U)} = \rho_U
$$
Taking $i = (U, V, U, a, \text{id}_U) \in \Omega$ for example, we find
that we have the desired compatibility. The uniqueness of $F$ follows
from the uniqueness of $E$ in the previous lemma (small detail omitted).
\end{proof}

\end{document}
